problem:  
Let \( n \ge 3 \), and let \( S^{n+1} \subset \mathbb{R}^{n+2} \) carry the round metric.  
Fix a compact Lie subgroup \( G \subset O(n+2) \) acting isometrically with cohomogeneity \( k \ge 2 \), and let \( Q = S^{n+1}/G \) denote the resulting compact Riemannian orbifold.  
Denote by \( \mathcal{M}_G(S^{n+1}) \) the moduli space (modulo isometries) of embedded \(G\)-invariant minimal hyperspheres in \( S^{n+1} \).

Each \(G\)-invariant minimal hypersphere \(\Sigma \in \mathcal{M}_G(S^{n+1})\) is the total preimage under the projection \(\pi:S^{n+1} \to Q\) of a smooth codimension-one suborbifold \(H_\Sigma \subset Q\) that is stationary for a natural weighted area functional taking into account the mean curvature of the orbits.  
The \(G\)-equivariant Jacobi operator on \(\Sigma\) corresponds to an elliptic operator on functions on \(H_\Sigma\) whose coefficients reflect local orbit geometry and curvature of \(Q\).

**Open Problem (Equivariant Index Blow-Up Near Singular Orbit Strata):**  
Does there exist a compact subgroup \( G \subset O(n+2) \) with cohomogeneity \( k \ge 2 \) and a sequence of embedded \( G \)-invariant minimal hyperspheres \(\Sigma_j = \pi^{-1}(H_{\Sigma_j}) \subset S^{n+1}\) satisfying:  
1. The projected hypersurfaces \( H_{\Sigma_j} \subset Q \) converge (e.g. as varifolds) to a limit suborbifold \(H_\infty\) whose support lies **in or arbitrarily close to** a singular orbit stratum of \(Q\);  
2. The areas \( \mathcal{A}(\Sigma_j) \) remain uniformly bounded; yet  
3. Their \(G\)-equivariant Morse indices \( \operatorname{ind}_G(\Sigma_j) \to \infty \)?  

Intuitively: Can degeneration of the orbit‑space geometry toward singular orbits cause unbounded spectral instability (Morse index growth) in invariant minimal hyperspheres, even when those hyperspheres do not experience geometric blow-up in area or curvature?  
Or does the structure of the orbifold metrics near the singular locus prevent such index blow‑up?

Why is it a "good" problem:  
This refined version is mathematically consistent and conceptually deep.  
It asks whether symmetry‑induced singularities in the orbit space \(Q=S^{n+1}/G\) can generate analytic instability for families of \(G\)-invariant minimal hyperspheres, bridging several important areas:

- **Equivariant geometric measure theory:** understanding compactness, varifold convergence, and degeneration under symmetry.
- **Spectral geometry on singular spaces:** since the equivariant Jacobi operator yields, on \(H_\Sigma\), an elliptic operator with singular coefficients that depend on the orbit‑type decomposition near singular strata.
- **Stability and index theory under symmetry constraints:** relating bounded geometric quantities (area) to unbounded analytic ones (index).

The question is accessible to state yet potentially very difficult to answer, requiring integration of techniques from analysis on orbifolds, equivariant minimal surface theory, and spectral asymptotics.  
It directly probes how geometric singularities of the quotient interact with analytic stability properties—a subtle and unexplored frontier in modern geometric analysis.  
Hence it fits the ideal of a *good mathematical problem*: simple in formulation, logically sound, cross‑disciplinary, and deeply illuminating of the structures that govern minimal hypersurfaces under symmetry.